\section{Research Questions}
\label{sec:research_questions}

This research investigates whether established context-aware misinformation detection models, specifically DANES~\citep{truica2024danes} and GETAE~\citep{GETAE2025}, can be improved by replacing their text-only content branch with a multimodal content branch. From the gap identified in the literature, three research questions arise concerning the performance and computational cost of this substitution, the relative contribution of the content and social-context branches, and the dataset adaptation requirements for evaluating the proposed MACE framework.

\subsection*{RQ1: Multimodal Substitution and its Cost}
\label{subsec:rq1}

\begin{quote}
\emph{How does replacing a text-only content branch with a multimodal
content branch affect the performance of context-aware inauthentic
social media content detection models, and what is the associated cost
in terms of inference time and computational resource usage?}
\end{quote}

\paragraph{What it addresses.}
RQ1 isolates a single architectural change in the DANES/GETAE ensemble:
the text-only content branch is substituted with a multimodal branch
that processes both textual and visual input, while the social context
branch is left unchanged. By holding the social signal constant, any
observed change in detection performance can be attributed directly to
the multimodal substitution rather than to confounded changes elsewhere
in the architecture. The question also explicitly addresses the
computational cost of this substitution (including additional
inference time, GPU memory consumption, and parameter count), so that
the tradeoff between accuracy gain and resource overhead can be
quantified within a single controlled experimental frame.

\paragraph{Significance.}
Modern misinformation increasingly relies on images, memes,
screenshots, and out-of-context visual material, which text-only
content branches cannot interpret. Existing context-aware ensembles
such as DANES and GETAE pre-date this shift and may therefore be
leaving meaningful detection accuracy unrealised. Establishing whether
multimodal substitution measurably improves their performance, and at
what computational cost, gives the field an empirical basis for
deciding when richer content encoders justify their additional
resource requirements, a question that is increasingly important as
detection systems are deployed at platform scale.

\subsection*{RQ2: Branch Contribution Decomposition}
\label{subsec:rq2}

\begin{quote}
\emph{How does the relative contribution of the multimodal content
branch and the social context branch affect overall detection
performance?}
\end{quote}

\paragraph{What it addresses.}
RQ2 shifts the focus from substitution to decomposition. Given a
multimodal content branch and a propagation-graph social context
branch, the question asks how much each branch contributes to overall
detection performance and whether their signals are complementary or
redundant. It is answered through a controlled ablation in which each
branch is run in isolation and in combination, allowing the marginal
contribution of each modality and signal source to be quantified
rather than reported only as a single combined accuracy figure.

\paragraph{Significance.}
Most published context-aware detection models report a single combined
performance number and do not separate the contribution of content
modelling from the contribution of social-context modelling. This
makes it difficult for the field to determine where future research
effort is best invested: in richer content encoders, in deeper
propagation modelling, or in more sophisticated fusion mechanisms.
A principled decomposition on the proposed backbone therefore has
direct design implications for the next generation of misinformation
detection systems and turns an otherwise opaque ensemble into a more
interpretable one.

\subsection*{RQ3: Dataset Adaptation and Generalisability}
\label{subsec:rq3}

\begin{quote}
\emph{What criteria and preprocessing steps are required to adapt an
existing multimodal social-media dataset for use with the DANES/GETAE
ensemble backbone, and how does dataset choice affect the
generalisability of results?}
\end{quote}

\paragraph{What it addresses.}
RQ1 and RQ2 both assume that suitable multimodal datasets are available
and that the framework behaves consistently across them. RQ3 examines
this assumption directly. DANES and GETAE were originally evaluated on
Twitter15 and Twitter16, which provide propagation graphs but without
multimodal content, while datasets such as
FakeNewsNet~\citep{shu2020fakenewsnet} and
MuMin~\citep{mumindataset} offer richer multimodal content but
carry known practical limitations including broken or missing image
URLs, sparse or delayed user engagement traces, and class imbalance
inherited from fact-checking sources. RQ3 will establish explicit
criteria (such as a minimum image-availability rate, graph-density
thresholds, and class-balance requirements) under which a dataset is
considered suitable for MACE evaluation, and will compare performance
across the adapted datasets to determine how strongly dataset
characteristics influence the conclusions drawn from RQ1 and RQ2.

\paragraph{Significance.}
Cross-dataset generalisation is a recurring concern in misinformation
detection research, where strong performance on one benchmark frequently
fails to transfer to another~\citep{mumindataset}. Without explicit
attention to dataset adaptation, the contribution of multimodal content
modelling cannot be cleanly separated from artefacts of dataset
construction. Establishing transparent adaptation criteria therefore
serves two purposes: it makes the MACE results reproducible and
comparable across datasets, and it offers a methodological contribution
to the wider field, where dataset suitability is often treated as an
unstated pre-condition rather than a research variable. For MACE, the simultaneous requirement for text, image, and network modalities across both branches substantially narrows the usable candidate pool from the broader landscape of available fake news datasets \citep{lv2025multimodal}; determining which datasets are genuinely suitable, and under what preprocessing conditions, is therefore a non-trivial research problem in its own right.

Together, RQ1, RQ2, and RQ3 form a coherent and complementary
investigation. RQ1 establishes whether multimodal substitution
measurably improves context-aware misinformation detection on the
DANES/GETAE backbone and quantifies the cost of doing so. RQ2 explains
\emph{why} by decomposing how content and social signals interact and
contribute to overall performance. RQ3 grounds both findings by
determining whether observed gains generalise across multimodal
datasets adapted under explicit, transparent criteria. Answering all
three questions is necessary to make a defensible and actionable
contribution to the context-aware fake news detection literature.
